\documentclass[10pt,a4paper]{amsart}
\usepackage[margin=19mm]{geometry}
\usepackage{amsmath,amssymb,mathtools}
\usepackage[scr=rsfs,cal=euler]{mathalfa}
\usepackage{xfrac}
\usepackage[unicode,hidelinks,bookmarks=false]{hyperref}
\newcommand{\ie}{i.e.\ }
\newcommand{\cf}{cf.\ }
\newcommand{\resp}{respectively\ }
\newcommand{\Subsection}{\subsection}
\providecommand{\nobreakdash}{\nobreak\hbox{-}}
\setlength{\emergencystretch}{2em}
\multlinegap0pt
\usepackage{amssymb}
\usepackage{mathtools}
\usepackage{booktabs}
\usepackage{algorithm}\usepackage{algpseudocode}
\usepackage[shortlabels]{enumitem}
\usepackage{caption}
\usepackage{cancel}
\usepackage{tikz}
\usepackage{multirow}
\newtheorem{theorem}{THEOREM}
\newcommand{\R}{\mathbb{R}}
\title{Deep Network Trainability via Persistent Subspace Orthogonality}
\author{Alex Massucco, Davide Murari, Carola-Bibiane Sch\"onlieb}
\date{Source: arXiv:2508.02882v2 (CC BY 4.0)}
\begin{document}
\maketitle

\noindent\emph{Source and licence.} Deep Network Trainability via Persistent Subspace Orthogonality. Version arXiv:2508.02882v2 (CC BY 4.0), distributed by arXiv under the Creative Commons Attribution 4.0 International licence, http://creativecommons.org/licenses/by/4.0/, as recorded in the arXiv licence metadata for this exact version and shown on the abstract page https://arxiv.org/abs/2508.02882v2. Full licence notice: \texttt{licenses/arxiv-2508.02882-CC-BY-4.0.txt}.\par\smallskip

\noindent\textit{Editorial scope.} Counted unit: the article's theorem thm:density (source0.tex lines 388--394). The article's complete original proof of it is delivered with it (source0.tex lines 977--1049, one \texttt{proof} environment of 73 lines, which the article places away from the statement). No mathematics is written, repaired or re-derived: the statement and the proof are the article's own lines, in the article's own notation, and no retained line is substituted or re-flowed.  Retained uncounted as the source-local context the proof needs: lines 378--380; lines 966--968.  Nothing was omitted from any retained span, so the package carries no omission record.  No retained line contains a citation, so the wrapper carries no bibliography and no bibliography entry.  No retained line contains a figure environment, an image-inclusion command or drawing code, so no image is delivered and the package declares no asset dependency.  Editorial formatting only, in addition: an AMS \texttt{amsart} preamble that restates the article's own declarations and macros used by the retained lines, namely the environments \texttt{theorem} on separate counters, together with the standard packages those lines need.  The retained environments print in this document as Theorem 1 in order of appearance, whereas the article prints the same items with the numbering of its own full document.  The complete native source of the article as distributed in the arXiv e-print for this exact version is bundled unchanged as \texttt{sources/182701/source0.tex}.
\begin{align}
&\mathcal{F}_\Omega(x; \widetilde{m}, \widetilde{q})=F(x) = \widetilde{m}(x)x + \widetilde{q}(x) + 1_{\Omega}(x)B^\top \sigma(Bx+b;\,\,1-\widetilde{m}(x),-1-\widetilde{m}(x)),\label{eq:pwcontinuous}
\end{align}

\begin{equation}\label{eq:pwl}
\sigma(x) = \max\{\sigma^1(x),\cdots,\sigma^k(x)\},
\end{equation}

\begin{theorem}\label{thm:density}
Let us consider a compact and connected set $\Omega\subset\R^n$, $B\in\mathcal{O}(n)$, $b\in\R^n$, and $\varepsilon>0$. Then, for any pair of continuous functions $m,q:\Omega\to\R$, there exists $N\in\mathbb{N}$, a finite collection of subsets $\{\Omega_i\subseteq\Omega|\,i=1,\cdots,N\}$ in $\mathcal{P}(\Omega)$, and a choice of $2N$ scalars $m_1,\cdots,m_N,q_1,\cdots,q_N\in\R$ such that
\[
\max_{x\in\Omega}\left\|\mathcal{F}_\Omega(x; m, q) - \mathcal{F}_\Omega(x; \widetilde{m}, \widetilde{q})\right\|_{\infty} < \varepsilon,
\]
where $\mathcal{F}_\Omega(x; \cdot, \cdot)$ is given by \eqref{eq:pwcontinuous} and $\widetilde{m}(x) = \sum_{i=1}^N 1_{\Omega_i}(x)m_i$, $\widetilde{q}(x) = \sum_{i=1}^N 1_{\Omega_i}(x)q_i$.
\end{theorem}

\begin{proof}[Proof of Theorem \ref{thm:density}]
Assume $\Omega$ is compact, and consider an activation function $\sigma$ as in \eqref{eq:pwl} with slopes limited to a set of two elements. %By the Heine-Cantor Theorem any continuous function $g:\Omega\to\R$ is uniformly continuous, and hence, for any $\varepsilon>0$, there exists a finite collection of subsets $\{\Omega_1,\cdots,\Omega_N\}$ satisfying $\mathcal{P}(\Omega)$ such that $\|g(x)-g(y)\|<\varepsilon$ for every $x,y\in\Omega_i$ and every $i=1,\cdots,N$. 
Fix $\varepsilon>0$ and two continuous functions $m,q:\Omega\to\mathbb{R}$. Call $M_{\Omega}:=\max_{x\in\Omega}\|x\|_2$ and consider a finite collection of sets $\{\Omega_1,\cdots,\Omega_N\}$ in $\mathcal{P}(\Omega)$ such that
\[
\begin{aligned}
\max_{x\in\Omega}|m(x) - \widetilde{m}(x)| &\leq \frac{\varepsilon}{4\sqrt{n}(M_{\Omega}\sqrt{n}+\|b\|_2)},\\
\max_{x\in\Omega}|q(x) - \widetilde{q}(x)| &\leq \frac{\varepsilon}{2},
\end{aligned}
\]
where $\widetilde{m},\widetilde{q}:\Omega\to\mathbb{R}$ are piecewise constant approximations of $m$ and $q$, which are constant on each of the sets $\Omega_1,\cdots,\Omega_N$.

%Thus,
%\[
%\max_{x\in\Omega}\left |g(x)- \sum_{i=1}^N 1_{\Omega_i}(x)g(x_i)\right | <\varepsilon,
%\]
%where $x_i\in \Omega_i$ is an arbitrary point. Up to taking a fine enough partition, we conclude that for any $\varepsilon>0$, and for any pair of continuous functions $m,q\in\mathcal{C}^0(\Omega,\R)$, there is a partition into open sets $\Omega_1,\cdots,\Omega_N$, and two piecewise constant functions $\widetilde{m},\widetilde{q}:\Omega\to \R$ defined as in \eqref{eq:pwcontinuous} giving
%\[
%\max_{x\in\Omega}\left\| m(x)x+q(x)-\widetilde{m}(x)x-\widetilde{q}(x) \right\|_{\infty} < \varepsilon.
%\]
Consider now the two functions $\mathcal{F}_{\Omega}(x;m,q)$ and $\mathcal{F}_{\Omega}(x;\widetilde{m},\widetilde{q})$
and write the difference
\[
\begin{split}
&\mathcal{F}_{\Omega}(x;m,q)-\mathcal{F}_\Omega(x;\widetilde{m},\widetilde{q})\\
& \qquad \qquad = (m(x)-\widetilde{m}(x))x + (q(x)-\widetilde{q}(x))\\
& \qquad \qquad +1_{\Omega}(x)B^\top\Bigl(\sigma(Bx+b;1-m(x),-1-m(x))\\
& \qquad \qquad - \sigma(Bx+b;1-\widetilde{m}(x),-1-\widetilde{m}(x))\Bigr).
\end{split}
\]
We now explicitly write the two functions whose difference needs to be controlled:
\[
\begin{split}
& \sigma(Bx+b;1-m(x),-1-m(x))= \max_{i=1,\cdots,k}\min_{j=1,\cdots,\ell_i}(a_{i,j}(x)(Bx+b) + b_{i,j}) \\
&\sigma(Bx+b;1-\widetilde{m}(x),-1-\widetilde{m}(x)) = \max_{i=1,\cdots,k}\min_{j=1,\cdots,\ell_i}(\widetilde{a}_{i,j}(x)(Bx+b) + {b}_{i,j}).
\end{split}
\]
Since the expression based on $\widetilde{m}$ and $\widetilde{q}$ can be constructed specifically for the target function at hand, we can assume that $a_{i,j}$ and $\widetilde{a}_{i,j}$ have agreeing expressions, meaning that if $a_{i,j}(x)=1-m(x)$, then $\widetilde{a}_{i,j}(x)=1-\widetilde{m}(x)$, for example. It follows that, for any $x\in\Omega$, one has
\[
\begin{split}
&\|\sigma(Bx+b;1-m(x),-1-m(x))-\sigma(Bx+b;1-\widetilde{m}(x),-1-\widetilde{m}(x))\|_\infty\\
&=\max_{r=1,\cdots,n} |\sigma((Bx+b)_r;1-m(x),-1-m(x))-\sigma((Bx+b)_r;1-\widetilde{m}(x),-1-\widetilde{m}(x))| \\
&\leq \max_{r=1,\cdots,n} \max_{i,j}|(a_{i,j}(x)-\widetilde{a}_{i,j}(x))(Bx+b)_r| = \max_{i,j} \max_{r=1,\cdots,n} |(a_{i,j}(x)-\widetilde{a}_{i,j}(x))(Bx+b)_r|\\
&= \max_{i,j} \|(a_{i,j}(x)-\widetilde{a}_{i,j}(x))(Bx+b)\|_\infty\leq \left(\max_{i,j} |a_{i,j}(x)-\widetilde{a}_{i,j}(x)|\right )\|Bx+b\|_\infty ,
\end{split}
\]
where $a_{i,j}(x)\in\{1-m(x),-1-m(x)\}$, $\widetilde{a}_{i,j}(x)\in\{1-\widetilde{m}(x),-1-\widetilde{m}(x)\}$, and $b_{i,j}\in\mathbb{R}$. This means that
\[
\begin{aligned}
\max_{i,j}|a_{i,j}(x)-\widetilde{a}_{i,j}(x)| & \leq |m(x)-\widetilde{m}(x)|\leq \frac{\varepsilon}{4\sqrt{n}(M_{\Omega}\sqrt{n}+\|b\|_2)}.
\end{aligned}
\]
Thus, it follows that
\[
\begin{aligned}
&\|\left(\sigma(Bx+b;1-m(x),-1-m(x)) -\sigma(Bx+b;1-\widetilde{m}(x),-1-\widetilde{m}(x))\right)\|_\infty\\
&\leq \sqrt{n} \frac{\varepsilon}{4\sqrt{n}(M_\Omega\sqrt{n}+\|b\|_2)}(M_\Omega\sqrt{n}+\|b\|_2) = \frac{\varepsilon}{4}.
\end{aligned}
\]

We then obtain the desired result
\[
\begin{aligned}
&\max_{x\in\Omega}\|\mathcal{F}_{\Omega}(x;m,q)-\mathcal{F}_\Omega(x;\widetilde{m},\widetilde{q})\|_\infty\leq \frac{\varepsilon}{4\sqrt{n}(M_\Omega\sqrt{n} + \|b\|_2)}M_\Omega + \frac{\varepsilon}{2} + \frac{\varepsilon}{4}\leq \frac{\varepsilon}{4} + \frac{\varepsilon}{2} + \frac{\varepsilon}{4} = \varepsilon
\end{aligned}
\]
as desired, where we used the estimate
\[
\begin{split}
\|(m(x)-\widetilde{m}(x))x\|_\infty &\leq \max_{x\in\Omega}|m(x)-\widetilde{m}(x)| \|x\|_2\leq \frac{\varepsilon}{4\sqrt{n}(M_\Omega\sqrt{n} + \|b\|_2)}M_\Omega\leq \frac{\varepsilon \sqrt{n}(M_\Omega\sqrt{n}+\|b\|_2)}{4\sqrt{n}(M_\Omega\sqrt{n} + \|b\|_2)} = \frac{\varepsilon}{4}.
\end{split}
\]
This allows us to conclude the proof.
\end{proof}

\end{document}
